% !TEX root = ../../cookbook.tex
\newpage
\recipe{Tarte au Vin Cuit}{ou ``Tarte au \textit{Raisinée}''.}{}{}

\subsection*{Ingrédients}
\begin{ingredients}
	\item 250 g de pâte sucrée (voir la recette à la page \pageref{bundner_nusstorte})
	\item 3 \oe ufs
	\item 1 jaune d'\oe uf
	\item 1.4 dl de crème double
	\item 1 dl de vin cuit vaudoise
\end{ingredients}

\medskip

% --- Preparation ---
Préparer la \textbf{pâte sucrée} et laisser-la reposer. 
Abaisser la pâte à \sfrac{1}{2} cm d'épaisseur et foncer une plaque de 20 cm de \diameter, beurrée et farinée. 
Faire cuire la pâte à blanc\footnote{Pour éviter que le fond de la pâte ne gonfle pendant la cuisson à blanc, le masquer par un disque de papier aluminium ou cuisson, et couvrir avec des pois chiches pour bien le tenir en place.}, 25 minutes environ, dans le four à 250°C. 
Laisser refroidir sans démouler. 

Préparer l'appareil à flan en mélangeant au fouet le \textbf{reste des ingrédients}\footnote{Selon le degré de réduction du vin cuit, il faudra peut-être en augmenter la quantité jusqu'à 1.2 dl.}. 
Remplisser le fond de pâte, transporter-la dans le four préchauffé à 180°C et achever de remplir à ras bords, une fois le gateau dans le four, en s'aidant d'une cuillère.

Laisser la porte du four entrouverte jusqu'à ce que le flan soit pris, une démi-heure environ, et vérifier son degré de cuisson en secouant légèrement le moule. 
Démouler à tiède sur une grille.


\vspace*{0.75 cm}

% --- Description ---
\begin{recipeblock}
	\noindent 
\end{recipeblock}